\documentclass[10pt,a4paper]{amsart}
\usepackage[margin=19mm]{geometry}
\usepackage{amsmath,amssymb,mathtools}
\usepackage[scr=rsfs,cal=euler]{mathalfa}
\usepackage{xfrac}
\usepackage[unicode,hidelinks,bookmarks=false]{hyperref}
\newcommand{\ie}{i.e.\ }
\newcommand{\cf}{cf.\ }
\newcommand{\resp}{respectively\ }
\newcommand{\Subsection}{\subsection}
\providecommand{\nobreakdash}{\nobreak\hbox{-}}
\setlength{\emergencystretch}{2em}
\multlinegap0pt
\usepackage[shortlabels]{enumitem}
\usepackage{mathtools}
\usepackage{dsfont}
\usepackage{amssymb}
\newtheorem{proposition}{Proposition}
\newtheorem{lemma}{Lemma}
\newcommand{\Buo}{\ensuremath{\mathsf{Buo}}}
\newcommand{\N}{\mathbb{N}}
\newcommand{\R}{\mathbb{R}}
\renewcommand{\ge}{\geqslant}
\renewcommand{\le}{\leqslant}
\title{Bourgain--uo sequential completeness in vector lattices}
\author{Tomasz Kania, Jaros{\l}aw Swaczyna}
\date{Source: arXiv:2512.14949v2 (CC BY 4.0)}
\begin{document}
\maketitle

\noindent\emph{Source and licence.} Bourgain--uo sequential completeness in vector lattices. Version arXiv:2512.14949v2 (CC BY 4.0), distributed by arXiv under the Creative Commons Attribution 4.0 International licence, http://creativecommons.org/licenses/by/4.0/, as recorded in the arXiv licence metadata for this exact version and shown on the abstract page https://arxiv.org/abs/2512.14949v2. Full licence notice: \texttt{licenses/arxiv-2512.14949-CC-BY-4.0.txt}.\par\smallskip

\noindent\textit{Editorial scope.} Counted unit: the article's proposition prop:Lipb-uniformly-discrete (source0.tex lines 809--813). The article's complete original proof of it is delivered with it (source0.tex lines 815--975, one \texttt{proof} environment of 161 lines). No mathematics is written, repaired or re-derived: the statement and the proof are the article's own lines, in the article's own notation, and no retained line is substituted or re-flowed.  Retained uncounted as the source-local context the proof needs: lines 182--188; lines 285--288; lines 707--720; lines 779--784.  Nothing was omitted from any retained span, so the package carries no omission record.  No retained line contains a citation, so the wrapper carries no bibliography and no bibliography entry.  No retained line contains a figure environment, an image-inclusion command or drawing code, so no image is delivered and the package declares no asset dependency.  Editorial formatting only, in addition: an AMS \texttt{amsart} preamble that restates the article's own declarations and macros used by the retained lines, namely the environments \texttt{proposition}, \texttt{lemma} on separate counters, together with the standard packages those lines need.  The retained environments print in this document as Proposition 1, Lemma 1 in order of appearance, whereas the article prints the same items with the numbering of its own full document.  The complete native source of the article as distributed in the arXiv e-print for this exact version is bundled unchanged as \texttt{sources/191133/source0.tex}.
\begin{proposition}\label{prop:Buo=order}
For a sequence $(x_n)$ in a vector lattice $E$ and $x\in E$ the following are equivalent:
\begin{enumerate}[label=(\roman*)]
\item $x_n\xrightarrow{\Buo} x$;
\item $x_n\xrightarrow{\mathrm{o}} x$.
\end{enumerate}
\end{proposition}

\begin{lemma}\label{lem:monotone_order_limit}
Let $E$ be a vector lattice and let $(x_n)$ be an increasing sequence in $E$.
If $x_n\xrightarrow{\mathrm{o}}x$, then $x=\sup_n x_n$ in $E$.
\end{lemma}

\begin{lemma}\label{lem:UC-approx-by-Lip}
Let $(X,d)$ be a metric space and let $g:X\to\R$ be bounded and uniformly continuous.
For $n\in\N$, define the \emph{inf-convolution}
\[
g_n(x):=\inf_{y\in X}\bigl(g(y)+n\,d(x,y)\bigr)\qquad(x\in X).
\]
Then:
\begin{enumerate}[\upshape(i)]
\item each $g_n$ is bounded and $n$-Lipschitz;
\item $g_n\le g$ pointwise;
\item $\|g_n-g\|_\infty\to0$ as $n\to\infty$;
\item $(g_n)$ is increasing: if $m\ge n$, then $g_n\le g_m$ pointwise.
\end{enumerate}
\end{lemma}

\begin{lemma}\label{lem:many-close-pairs}
Let $(X,d)$ be a metric space with $\delta(X)=0$ and $d(x)>0$ for all $x\in X$
(i.e.,\ $X$ is discrete but not uniformly discrete). Then for every $\varepsilon>0$
and every finite set $F\subset X$, there exist distinct points $a,b\in X\setminus F$
with $d(a,b)<\varepsilon$.
\end{lemma}

\begin{proposition}\label{prop:Lipb-uniformly-discrete}
Let $(X,d)$ be a metric space and let $\mathrm{Lip}_b(X)$ be the vector lattice of bounded
Lipschitz functions $X\to\R$, ordered pointwise. Then $\mathrm{Lip}_b(X)$ is sequentially
\Buo-complete if and only if $\delta(X)>0$.
\end{proposition}

\begin{proof}
\noindent For necessity, assume $\delta(X)=0$. We construct a \Buo-Cauchy sequence in
$\mathrm{Lip}_b(X)$ with no order limit.

\smallskip\noindent
We first produce a closed set $A\subset X$ and a sequence $(b_n)\subset X\setminus A$ such that
\[
t_n:=\mathrm{dist}(b_n,A):=\inf_{a\in A}d(b_n,a)\in(0,1)\quad\text{and}\quad t_n\downarrow0.
\]

\smallskip
\emph{Case A: there exists $x_0\in X$ with $d(x_0)=0$.}
Then $x_0$ is non-isolated, so there exist $b_n\neq x_0$ with $d(b_n,x_0)\to0$.
Let $A:=\{x_0\}$. Then $t_n=\mathrm{dist}(b_n,A)=d(b_n,x_0)\downarrow0$; passing to a tail, assume $t_n<1$ for all $n$.

\smallskip
\emph{Case B: $d(x)>0$ for all $x\in X$.}
Then $X$ is discrete, and since $\delta(X)=0$ it is not uniformly discrete.
We build inductively pairs $(a_n,b_n)$ of distinct points such that
\[
d(a_n,b_n)<\frac1n
\quad\text{and}\quad
\{a_1,\dots,a_n\}\cap\{b_1,\dots,b_n\}=\varnothing.
\]
Choose any distinct $a_1,b_1$ with $d(a_1,b_1)<1$.
Assuming $(a_1,b_1),\dots,(a_{n-1},b_{n-1})$ chosen, let
\[
F:=\{a_1,\dots,a_{n-1},b_1,\dots,b_{n-1}\}.
\]
By Lemma~\ref{lem:many-close-pairs}, there exist distinct $a_n,b_n\in X\setminus F$ with $d(a_n,b_n)<1/n$.
Set $A:=\{a_n:n\in\N\}$. Then $b_n\notin A$ and
\[
t_n=\mathrm{dist}(b_n,A)\le d(b_n,a_n)<\frac1n<1,
\]
so $t_n\downarrow0$. Moreover $t_n>0$ since $b_n\notin A$ and $d(b_n)>0$.

\smallskip\noindent
In both cases, define the bounded uniformly continuous function
\[
g(x):=\sqrt{\mathrm{dist}(x,A)}\wedge 1\qquad(x\in X).
\]
(Here $x\mapsto \mathrm{dist}(x,A)$ is $1$-Lipschitz and $t\mapsto \sqrt{t}\wedge 1$
is uniformly continuous on $[0,\infty)$.)

\smallskip\noindent
\emph{$g$ is not Lipschitz.}
For each $n$, choose $a'_n\in A$ such that $d(b_n,a'_n)<2t_n$ (possible by the definition of $t_n$).
Then $g(a'_n)=0$ and $g(b_n)=\sqrt{t_n}$, hence
\[
\frac{|g(b_n)-g(a'_n)|}{d(b_n,a'_n)}
>\frac{\sqrt{t_n}}{2t_n}
=\frac{1}{2\sqrt{t_n}}
\longrightarrow\infty,
\]
so $g\notin\mathrm{Lip}_b(X)$.

\smallskip\noindent
\emph{Approximation by bounded Lipschitz functions.}
By Lemma~\ref{lem:UC-approx-by-Lip}, the inf-convolutions
\[
g_n(x):=\inf_{y\in X}\bigl(g(y)+n\,d(x,y)\bigr)
\]
belong to $\mathrm{Lip}_b(X)$ and satisfy $\|g_n-g\|_\infty\to0$.

\smallskip\noindent
\emph{$(g_n)$ is \Buo-Cauchy.}
Let $(n_k)$ be strictly increasing. Then
\[
\|g_{n_{k+1}}-g_{n_k}\|_\infty
\le \|g_{n_{k+1}}-g\|_\infty+\|g-g_{n_k}\|_\infty
\longrightarrow0.
\]
Define $\varepsilon_m:=\sup_{k\ge m}\|g_{n_{k+1}}-g_{n_k}\|_\infty$. Then $\varepsilon_m\downarrow0$, and for all $k\ge m$,
\[
|g_{n_{k+1}}-g_{n_k}|\le \varepsilon_m\,\mathds1_X,
\]
where $\mathds1_X$ is the constant $1$ function. Since $\varepsilon_m\mathds1_X\downarrow0$ in $\mathrm{Lip}_b(X)$,
it follows that $g_{n_{k+1}}-g_{n_k}\xrightarrow{\mathrm{o}}0$.
Thus $(g_n)$ is \Buo-Cauchy.

\smallskip\noindent
\emph{$(g_n)$ has no order limit in $\mathrm{Lip}_b(X)$.}
Suppose towards a contradiction that $g_n\xrightarrow{\mathrm{o}}h$ in
$\mathrm{Lip}_b(X)$. By Lemma~\ref{lem:UC-approx-by-Lip}\textup{(iv)} the
sequence $(g_n)$ is increasing, hence Lemma~\ref{lem:monotone_order_limit}
yields that $h$ is the supremum of $\{g_n:n\in\N\}$ in
$\mathrm{Lip}_b(X)$.

For each $m\in\N$, choose $N_m\in\N$ such that
$\|g-g_{N_m}\|_\infty<\tfrac1m$, and define
\[
u_m:=g_{N_m}+\frac1m\,\mathds{1}_X\in \mathrm{Lip}_b(X).
\]
Since $g_n\le g$ for all $n$ and $g\le u_m$ pointwise, each $u_m$ is an
upper bound of the set $\{g_n:n\in\N\}$ in $\mathrm{Lip}_b(X)$. By minimality
of the supremum, we therefore have
\[
h\le u_m\qquad(m\in\N).
\]
Fix $x\in X$. Since $g_{N_m}(x)\to g(x)$ and $m^{-1}\to 0$, it follows that
$u_m(x)\to g(x)$, hence $h(x)\le g(x)$. On the other hand, $h$ is itself an
upper bound of $\{g_n:n\in\N\}$, while $g_n(x)\uparrow g(x)$ pointwise, so
$h(x)\ge g(x)$. Thus $h(x)=g(x)$ for all $x\in X$, i.e.\ $h=g$.
This contradicts $g\notin \mathrm{Lip}_b(X)$.
Therefore $\mathrm{Lip}_b(X)$ is not sequentially \Buo-complete when
$\delta(X)=0$.

\medskip
\noindent For sufficiency, assume $\delta:=\delta(X)>0$.
Let $f:X\to\R$ be bounded. For distinct $x,y\in X$ we have $d(x,y)\ge\delta$, hence
\[
\frac{|f(x)-f(y)|}{d(x,y)}\le \frac{2\|f\|_\infty}{\delta}.
\]
Thus every bounded function is Lipschitz, and consequently
\[
\mathrm{Lip}_b(X)=\ell_\infty(X)
\]
as vector lattices (both are all bounded functions on $X$ with the pointwise order).

It therefore suffices to show that $\ell_\infty(X)$ is sequentially \Buo-complete.
Let $(f_n)\subset\ell_\infty(X)$ be \Buo-Cauchy.

\smallskip\noindent
\emph{Pointwise convergence.}
Fix $x\in X$. For any strictly increasing $(n_k)$, the differences
$f_{n_{k+1}}-f_{n_k}$ converge to $0$ in order in $\ell_\infty(X)$, hence in particular
\[
f_{n_{k+1}}(x)-f_{n_k}(x)\to0\quad\text{in }\R.
\]
Thus $(f_n(x))_n$ is a Cauchy sequence in $\R$, hence convergent.
Define $f(x):=\lim_{n\to\infty}f_n(x)$.

\smallskip\noindent
\emph{Uniform boundedness.}
We claim $\sup_n\|f_n\|_\infty<\infty$.
Suppose not. Choose a strictly increasing $(n_k)$ with
\[
\|f_{n_{k+1}}\|_\infty\ge \|f_{n_k}\|_\infty+(k+1)\qquad(k\in\N).
\]
Set $d_k:=f_{n_{k+1}}-f_{n_k}$. Then
\[
\|d_k\|_\infty\ge \|f_{n_{k+1}}\|_\infty-\|f_{n_k}\|_\infty\ge k+1,
\]
so $(d_k)$ is not norm bounded.
However $(d_k)$ converges to $0$ in order by \Buo-Cauchyness, hence it is order bounded:
there exists $0\le y\in\ell_\infty(X)$ with $|d_k|\le y$ for all $k$.
Then $\|d_k\|_\infty\le\|y\|_\infty$ for all $k$, a contradiction.
Thus $M:=\sup_n\|f_n\|_\infty<\infty$, and in particular $|f(x)|\le M$ for all $x$, so $f\in\ell_\infty(X)$.

\smallskip\noindent
\emph{Order convergence.}
For $m\in\N$ define
\[
z_m(x):=\sup_{n\ge m}|f_n(x)-f(x)|,\qquad x\in X.
\]
Then $0\le z_{m+1}\le z_m$ pointwise, and $z_m(x)\downarrow0$ for each $x$ since $f_n(x)\to f(x)$.
Also $z_m(x)\le 2M$ for all $x$, hence $z_m\in\ell_\infty(X)$.
Finally, for every $n\ge m$ we have $|f_n-f|\le z_m$, so $f_n\xrightarrow{\mathrm{o}}f$ in $\ell_\infty(X)$.
By Proposition~\ref{prop:Buo=order}, this is equivalent to $f_n\xrightarrow{\Buo}f$.
Therefore $\ell_\infty(X)$ (and hence $\mathrm{Lip}_b(X)$) is sequentially \Buo-complete.
\end{proof}

\end{document}
